% Options for packages loaded elsewhere
% Options for packages loaded elsewhere
\PassOptionsToPackage{unicode}{hyperref}
\PassOptionsToPackage{hyphens}{url}
%
\documentclass[
  letterpaper,
  12pt,
  oneside]{book}
\usepackage{xcolor}
\usepackage{amsmath,amssymb}
\setcounter{secnumdepth}{5}
\usepackage{iftex}
\ifPDFTeX
  \usepackage[T1]{fontenc}
  \usepackage[utf8]{inputenc}
  \usepackage{textcomp} % provide euro and other symbols
\else % if luatex or xetex
  \usepackage{unicode-math} % this also loads fontspec
  \defaultfontfeatures{Scale=MatchLowercase}
  \defaultfontfeatures[\rmfamily]{Ligatures=TeX,Scale=1}
\fi
\usepackage{lmodern}
\ifPDFTeX\else
  % xetex/luatex font selection
\fi
% Use upquote if available, for straight quotes in verbatim environments
\IfFileExists{upquote.sty}{\usepackage{upquote}}{}
\IfFileExists{microtype.sty}{% use microtype if available
  \usepackage[]{microtype}
  \UseMicrotypeSet[protrusion]{basicmath} % disable protrusion for tt fonts
}{}
\usepackage{setspace}
\makeatletter
\@ifundefined{KOMAClassName}{% if non-KOMA class
  \IfFileExists{parskip.sty}{%
    \usepackage{parskip}
  }{% else
    \setlength{\parindent}{0pt}
    \setlength{\parskip}{6pt plus 2pt minus 1pt}}
}{% if KOMA class
  \KOMAoptions{parskip=half}}
\makeatother
% Make \paragraph and \subparagraph free-standing
\makeatletter
\ifx\paragraph\undefined\else
  \let\oldparagraph\paragraph
  \renewcommand{\paragraph}{
    \@ifstar
      \xxxParagraphStar
      \xxxParagraphNoStar
  }
  \newcommand{\xxxParagraphStar}[1]{\oldparagraph*{#1}\mbox{}}
  \newcommand{\xxxParagraphNoStar}[1]{\oldparagraph{#1}\mbox{}}
\fi
\ifx\subparagraph\undefined\else
  \let\oldsubparagraph\subparagraph
  \renewcommand{\subparagraph}{
    \@ifstar
      \xxxSubParagraphStar
      \xxxSubParagraphNoStar
  }
  \newcommand{\xxxSubParagraphStar}[1]{\oldsubparagraph*{#1}\mbox{}}
  \newcommand{\xxxSubParagraphNoStar}[1]{\oldsubparagraph{#1}\mbox{}}
\fi
\makeatother


\usepackage{longtable,booktabs,array}
\usepackage{calc} % for calculating minipage widths
% Correct order of tables after \paragraph or \subparagraph
\usepackage{etoolbox}
\makeatletter
\patchcmd\longtable{\par}{\if@noskipsec\mbox{}\fi\par}{}{}
\makeatother
% Allow footnotes in longtable head/foot
\IfFileExists{footnotehyper.sty}{\usepackage{footnotehyper}}{\usepackage{footnote}}
\makesavenoteenv{longtable}
\usepackage{graphicx}
\makeatletter
\newsavebox\pandoc@box
\newcommand*\pandocbounded[1]{% scales image to fit in text height/width
  \sbox\pandoc@box{#1}%
  \Gscale@div\@tempa{\textheight}{\dimexpr\ht\pandoc@box+\dp\pandoc@box\relax}%
  \Gscale@div\@tempb{\linewidth}{\wd\pandoc@box}%
  \ifdim\@tempb\p@<\@tempa\p@\let\@tempa\@tempb\fi% select the smaller of both
  \ifdim\@tempa\p@<\p@\scalebox{\@tempa}{\usebox\pandoc@box}%
  \else\usebox{\pandoc@box}%
  \fi%
}
% Set default figure placement to htbp
\def\fps@figure{htbp}
\makeatother





\setlength{\emergencystretch}{3em} % prevent overfull lines

\providecommand{\tightlist}{%
  \setlength{\itemsep}{0pt}\setlength{\parskip}{0pt}}



 


% McGill thesis preamble
% Adapted from the McGill thesis LaTeX template
% https://github.com/juliengs/mcgill-thesis-template

% Big gray chapter numbers
\usepackage{fix-cm}
\usepackage{xcolor}
\usepackage[bf,rm,medium,compact]{titlesec}
\definecolor{gray75}{gray}{0.75}
\newcommand{\hsp}{\hspace{20pt}}
\titleformat{\chapter}[display]
  {\filleft\Huge\bfseries}
  {\fontsize{100}{100}\selectfont\textcolor{gray75}\thechapter}
  {1ex}{}[]

% Format for unnumbered chapters (frontmatter, backmatter)
\titleformat{name=\chapter,numberless}
  {\Huge\bfseries}
  {}
  {0pt}{}

% Suppress default title page (we use _templates/before-body.tex)
\renewcommand{\maketitle}{}

% McGill guidelines link command (used in placeholder text)
\newcommand{\mcgillguidelines}{%
  \href{https://www.mcgill.ca/gps/thesis/thesis-guidelines/preparation}%
  {Official McGill Guidelines}: }
\makeatletter
\@ifpackageloaded{bookmark}{}{\usepackage{bookmark}}
\makeatother
\makeatletter
\@ifpackageloaded{caption}{}{\usepackage{caption}}
\AtBeginDocument{%
\ifdefined\contentsname
  \renewcommand*\contentsname{Table of contents}
\else
  \newcommand\contentsname{Table of contents}
\fi
\ifdefined\listfigurename
  \renewcommand*\listfigurename{List of Figures}
\else
  \newcommand\listfigurename{List of Figures}
\fi
\ifdefined\listtablename
  \renewcommand*\listtablename{List of Tables}
\else
  \newcommand\listtablename{List of Tables}
\fi
\ifdefined\figurename
  \renewcommand*\figurename{Figure}
\else
  \newcommand\figurename{Figure}
\fi
\ifdefined\tablename
  \renewcommand*\tablename{Table}
\else
  \newcommand\tablename{Table}
\fi
}
\@ifpackageloaded{float}{}{\usepackage{float}}
\floatstyle{ruled}
\@ifundefined{c@chapter}{\newfloat{codelisting}{h}{lop}}{\newfloat{codelisting}{h}{lop}[chapter]}
\floatname{codelisting}{Listing}
\newcommand*\listoflistings{\listof{codelisting}{List of Listings}}
\makeatother
\makeatletter
\makeatother
\makeatletter
\@ifpackageloaded{caption}{}{\usepackage{caption}}
\@ifpackageloaded{subcaption}{}{\usepackage{subcaption}}
\makeatother
\usepackage{bookmark}
\IfFileExists{xurl.sty}{\usepackage{xurl}}{} % add URL line breaks if available
\urlstyle{same}
\hypersetup{
  pdftitle={Title of Your Thesis},
  pdfauthor={Santiago Fernandez},
  hidelinks,
  pdfcreator={LaTeX via pandoc}}


% Intentionally empty: suppresses default Quarto/Pandoc title page.
% The McGill title page is rendered via _templates/before-body.tex.
\begin{document}
\frontmatter
\maketitle

% McGill title page
% Conforms to: https://www.mcgill.ca/gps/thesis/thesis-guidelines/preparation
\thispagestyle{empty}
\begin{center}

\vspace*{0.5cm}

{\bfseries\LARGE Title of Your}
\vspace{0.15cm}

{\bfseries\LARGE Thesis Goes Here}
\vspace{1.8cm}

{\large Santiago Fernandez}
\\Doctor of Philosophy

\vspace{1cm}
\includegraphics[width=.6\linewidth]{images/mcgill_sig_red.pdf}

\vspace{0.5cm}
Department of Neurology and Neurosurgery\\
McGill University\\
Montreal, Quebec, Canada\\

\vspace{1.5cm}
\today

\vspace{1.0cm}
\noindent
A thesis submitted to McGill University in partial\\
fulfillment of the requirements of the degree of\\
Doctor of Philosophy

\vspace{1.0cm}
{\small \copyright\ Santiago Fernandez, \the\year{}}

\end{center}
\clearpage


\setstretch{1.5}
\mainmatter
\frontmatter

\bookmarksetup{startatroot}

\chapter*{Abstract}\label{abstract}
\addcontentsline{toc}{chapter}{Abstract}

\markboth{Abstract}{Abstract}

\href{https://www.mcgill.ca/gps/thesis/thesis-guidelines/preparation}{Official
McGill Guidelines}: If the language of the thesis is neither English nor
French (only allowed for specific language Units) then a third abstract
in the language of the thesis is required.

Abstracts in English and French are mandatory and must be text only,
i.e. no images, special characters (apart from the West European
character set excluding the ``OE'' and ``oe''), chemical or mathematical
formulae, or special formatting (e.g.~lists, tables). Abstracts have a
maximum limit of 4,000 characters.

\bookmarksetup{startatroot}

\chapter*{Abrégé}\label{abruxe9guxe9}
\addcontentsline{toc}{chapter}{Abrégé}

\markboth{Abrégé}{Abrégé}

\href{https://www.mcgill.ca/gps/thesis/thesis-guidelines/preparation}{Official
McGill Guidelines}: La même chose en français.

\bookmarksetup{startatroot}

\chapter*{Contribution of Authors}\label{contribution-of-authors}
\addcontentsline{toc}{chapter}{Contribution of Authors}

\markboth{Contribution of Authors}{Contribution of Authors}

\href{https://www.mcgill.ca/gps/thesis/thesis-guidelines/preparation}{Official
McGill Guidelines}: A doctoral thesis must clearly state the elements of
the thesis that are considered original scholarship and distinct
contributions to knowledge.

\begin{itemize}
\tightlist
\item
  Contributions of the student to each chapter must be explicitly
  stated.
\item
  Contributions of any co-authors to each chapter must be explicitly
  stated.
\end{itemize}

\bookmarksetup{startatroot}

\chapter*{Acknowledgements}\label{acknowledgements}
\addcontentsline{toc}{chapter}{Acknowledgements}

\markboth{Acknowledgements}{Acknowledgements}

\href{https://www.mcgill.ca/gps/thesis/thesis-guidelines/preparation}{Official
McGill Guidelines}: Among other acknowledgements, the student is
required to declare the extent to which assistance (paid or unpaid) has
been given by members of staff, fellow students, research assistants,
technicians, or others in the collection of materials and data, the
design and construction of apparatus, the performance of experiments,
the analysis of data, and the preparation of the thesis (including
editorial help).

\begin{itemize}
\tightlist
\item
  In addition, it is appropriate to recognize the supervision and advice
  given by the thesis supervisor(s) and advisors.
\end{itemize}

\tableofcontents
\listoffigures
\listoftables
\clearpage
\mainmatter

\bookmarksetup{startatroot}

\chapter{Introduction}\label{introduction}

\href{https://www.mcgill.ca/gps/thesis/thesis-guidelines/preparation}{Official
McGill Guidelines}: Clearly state the rationale and objectives of the
research.

\bookmarksetup{startatroot}

\chapter{Literature Review}\label{literature-review}

\href{https://www.mcgill.ca/gps/thesis/thesis-guidelines/preparation}{Official
McGill Guidelines}: The comprehensive review of the literature must
sufficiently demonstrate the student's knowledge of and expertise in
their research areas and should be broad enough to apply to each
research question in the thesis. The review of the literature can
additionally include various types of content, such as:

\begin{itemize}
\tightlist
\item
  A review providing a reader who is relatively less familiar with the
  research topic (e.g., an internal/external member of an oral defence
  committee with adjacent but not direct expertise) an introduction to
  the general domain.
\item
  An explanation of the overall rationale for how and why the subsequent
  studies were conducted. For example, the literature underlying the
  research questions must be sufficiently discussed.
\item
  A review of fundamental theories underlying the subsequently presented
  work, or to explain why certain approaches were not taken in the
  study(ies) presented.
\end{itemize}

The literature review must be in line with disciplinary expectations.
The review can be incorporated in the Introduction chapter, addressed in
a standalone chapter, or distributed across multiple chapters.

\bookmarksetup{startatroot}

\chapter{Discussion}\label{discussion}

\href{https://www.mcgill.ca/gps/thesis/thesis-guidelines/preparation}{Official
McGill Guidelines}: The discussion of findings must be in line with
disciplinary expectations. A comprehensive discussion is expected to be
a minimum of 10 pages, double-spaced for doctoral students and a minimum
of 5 pages, double-spaced for Master's students (including figures,
images, and tables). It pertains to the entirety of a thesis. The
discussion of findings should provide a final, overarching summary of
study themes, limitations, and future directions.

In the case of a manuscript-based thesis, the comprehensive discussion
should encompass all of the chapters of the thesis and should not be a
repetition of the individual chapters. This section can be used to
address issues not sufficiently covered in the preceding chapters or
papers (e.g., critiques raised by reviewers that could not be
incorporated into published works, or reintroducing discussion arguments
removed from published papers upon reviewer request). This section can
also be used to elaborate on the practical/applied aspects of published
findings in a manner that is more accessible to less expert readers.

\bookmarksetup{startatroot}

\chapter{Conclusions}\label{conclusions}

\href{https://www.mcgill.ca/gps/thesis/thesis-guidelines/preparation}{Official
McGill Guidelines}: Clearly state how the objectives of the research
were met and discuss implications of findings.

\bookmarksetup{startatroot}

\chapter*{Publications}\label{publications}
\addcontentsline{toc}{chapter}{Publications}

\markboth{Publications}{Publications}

This part is optional, but it gives a nice touch to list all the
publications (official venues down to poster sessions) throughout your
PhD.

\bookmarksetup{startatroot}

\chapter*{Acronyms}\label{acronyms}
\addcontentsline{toc}{chapter}{Acronyms}

\markboth{Acronyms}{Acronyms}

\begin{itemize}
\tightlist
\item
  \textbf{HVR}: Hippocampal-to-Ventricle Ratio
\item
  \textbf{QC}: Quality Control
\item
  \textbf{CVR}: Cerebrovascular Reactivity
\item
  \textbf{LGCM}: Latent Growth Curve Model
\end{itemize}

\cleardoublepage
\phantomsection
\addcontentsline{toc}{part}{Appendices}
\appendix

\chapter{Supplementary Material for Chapter 3}\label{sec-appendix-ch3}

Supplementary tables and figures for the HVR Validation study (Chapter
3) will be placed here when the chapter is adapted.

\chapter{Supplementary Material for Chapter 4}\label{sec-appendix-ch4}

Supplementary tables and figures for the sex differences study (Chapter
4) will be placed here when the chapter is adapted.


\backmatter


\end{document}
